\documentclass[11pt,a4paper]{article}
\usepackage[margin=22.1mm]{geometry}
\usepackage[T1]{fontenc}
\usepackage{lmodern}
\usepackage{booktabs,tabularx}
\usepackage[hidelinks]{hyperref}

\hypersetup{
  pdftitle={Cardinality Feedback in PostgreSQL},
  pdfauthor={Prabhav Khandelwal, Sanchita Chaurasia, Saksham Saharia}
}
\urlstyle{same}
\setlength{\parindent}{0pt}
\setlength{\parskip}{5pt}

\begin{document}

\begin{center}
{\Large\bfseries Cardinality Feedback in PostgreSQL}\\[6pt]
Prabhav Khandelwal (23B0988)\\
Sanchita Chaurasia (23B1048)\\
Saksham Saharia (23B1078)\\[4pt]
CS631 Database Systems Project
\end{center}

\textbf{Project repository}\quad
\url{https://github.com/sanchitachaurasia/CS-631}

\section*{What we built}
Before PostgreSQL runs a query, it estimates how many rows each scan and join
will return. The planner uses those estimates to compare possible plans. For
example, a nested loop might be cheap when a join has only a few rows to
process. The same plan can do a lot of extra work if that estimate is too low.
Estimates can be especially difficult when filters on different columns or
tables are related.

We added a feedback loop to PostgreSQL 16.13. The executor measures how many
rows a completed scan or inner join produced. A later query can use that
measurement if it needs the same rows from the same view of the data. The
saved value is a row count. The later query still reads its data and executes
its chosen plan. The idea of learning from execution measurements comes from
LEO~\cite{leo}.

The initial prototype collected counts and reused them in a restricted
read-only transaction. We extended it so the count from a standalone table
scan can help plan a join. We also made it possible to reuse the count from a
two-table join when that pair appears inside a larger join. When the database
snapshot allows it, these observations can be shared across connections.

\section*{How it fits into PostgreSQL}
\textbf{At planning time.} In \texttt{costsize.c}, PostgreSQL computes its
usual row estimate for a table scan or an inner join. The new
\texttt{CardFeedbackEstimate} hook looks up a matching count when application
is enabled. On a hit it replaces the relation's estimate before paths are
costed. A count of zero becomes an estimate of one, following PostgreSQL's
minimum. PostgreSQL then uses its ordinary cost model and plan search.

\textbf{From the plan to the executor.} For an eligible scan or join,
\texttt{createplan.c} carries the lookup key and the estimate from before
replacement into the plan. Paths that take parameters from an outer relation
are skipped. Such a path might read different rows each time it is called, so a
count from one invocation would describe only that input.

\textbf{After execution.} With collection enabled, the executor records
complete observations in a shared store for later planning.

Both feedback settings default to off.

\begin{tabular}{@{}ll@{}}
\texttt{enable\_card\_feedback\_collect} & Records complete counts. \\
\texttt{enable\_card\_feedback\_apply} & Uses saved counts during planning. \\
\end{tabular}

The modified source is in \texttt{postgres/}. The first commit is the stock
baseline.\\
The patch in \texttt{project/postgres.patch} shows the same changes against
that baseline.

\newpage

\section*{Collecting and storing counts}
In \texttt{execMain.c}, the collection setting enables PostgreSQL's existing
row counters for a running \texttt{SELECT}. A plain \texttt{EXPLAIN} displays
a plan without running it, so it supplies no observations. The executor records
whether it ran forward to the end in one unrestricted call. It skips collection
after bounded reads, rewinds, or repeated executor calls.

At shutdown, \texttt{CardFeedbackCollect} checks the plan while its counters
are still available. The collector starts at the completed plan output and
follows only inputs that were fully read. Even then, a node must have run once,
because a counter accumulated over several loops can combine different
inputs. Some nodes require extra care.

\begin{itemize}
\item A nested loop reads its outer input, but it may stop an inner scan after
  the first match. That inner count can be smaller than the full scan count.
\item A serial hash build reads its whole build input. The collector follows
  the probe input only when the hash table is nonempty and uses one batch.
\item A merge join can stop reading either side early, and \texttt{Memoize}
  can return saved output from earlier calls. The collector does not follow
  those inputs.
\end{itemize}

\texttt{card\_feedback.c} keeps up to 1,024 entries in shared memory. Each
key has a 16 KiB identity buffer and a ring containing its 16 latest complete
observations. Readers can share the lock; updates and least-recently-updated
eviction use the exclusive lock. Each sample keeps the original estimate,
actual count, timestamp and sequence number.

The C code summarizes the ring with median and 90th-percentile Q-error,
log-scale variability, outlier count, drift between older and newer samples,
and a row-count trend. Confidence rises as evidence accumulates and falls
with instability, outliers, drift and age; its age component halves every
seven days. During planning, PostgreSQL blends its normal estimate with the
latest measured count according to this confidence.

The extension exposes these summaries through
\texttt{card\_feedback\_entries()} and the chronological samples through
\texttt{card\_feedback\_history()}. History is copied while holding the shared
lock and formatted after releasing it. \texttt{card\_feedback\_prune(seconds)}
removes entries whose last observation is older than the requested age.

\section*{Matching the same work}
A key represents the rows that a scan or join should produce. It includes the
effective restrictions from the planner, including filters that PostgreSQL
infers from other conditions in the query.

\begin{tabularx}{\linewidth}{@{}p{0.20\linewidth}X@{}}
\toprule
\textbf{Key} & \textbf{What is matched} \\
\midrule
Scan & The table, its storage locator and its effective filters. A change to
a column, operator or literal changes the key. An alias or a different
ordering of the top-level \texttt{AND} terms can still match. \\
\addlinespace
Two-table join & The scan keys for two distinct tables and the effective
column equalities between them. The key puts the tables and equality
operands in a consistent order. \\
\addlinespace
Query-specific join & The normalized join tree, the tables it refers to and
the subset of tables at this join. This covers other eligible joins within
the same query context. \\
\bottomrule
\end{tabularx}

\newpage

\section*{When reuse is allowed}
The query must run as a direct \texttt{SELECT} in an explicit read-only
\texttt{REPEATABLE READ} transaction with no assigned transaction ID and with
parallel query disabled. It must read permanent heap tables without
inheritance or partitioning. Prepared statements, subqueries, row locking,
row-level security, \texttt{LIMIT}, and \texttt{OFFSET} fall outside the reuse
rules. When a check fails or the key is missing, PostgreSQL keeps its ordinary
estimate.

\subsection*{Supported filters and joins}
\textbf{Table filters.}
Examples include \texttt{=}, \texttt{<>}, \texttt{<}, \texttt{<=},
\texttt{>}, and \texttt{>=} on ordinary columns with the built-in operators.
The prototype also handles \texttt{BETWEEN}, \texttt{IN} lists,
\texttt{= ANY} with an array, built-in \texttt{LIKE}, null and Boolean tests,
and combinations with \texttt{AND}, \texttt{OR}, and \texttt{NOT}.
The operator must use a built-in immutable function, which PostgreSQL marks
as returning the same result for the same inputs. Direct function calls
are limited to integer \texttt{abs} and casts between integer types.
Sequential, index, index-only, and bitmap heap scans can supply counts.
A scan count can carry over when the table appears in a join and its
effective filters still match.

\textbf{Inner joins.}
The prototype can collect complete output counts from supported nested-loop,
hash, and merge joins. A two-table count can transfer to a larger join
when it matches the effective filters on both inputs and has one or more
column equalities, such as \texttt{a.id = b.id}, that PostgreSQL recognizes
for merge joining two distinct tables. An inequality join such as
\texttt{a.id > b.id}, a self-join, or a \texttt{CROSS JOIN} can reuse its
join count in the same query context. Joins that remain outer joins are
outside the reuse rules.

A database snapshot describes which transactions' changes the query can
see. When a normal snapshot's lower and upper transaction ID bounds are
equal, the code adds the cutoff and its epoch to the key. Two connections
with the same cutoff can share a count under the supported conditions.
If an older transaction keeps those bounds apart, reuse stays within the
local transaction. A table rewrite changes the storage locator in the
key. A write elsewhere can move the cutoff even if the tables in the query
are unchanged.

\section*{An example with IMDb data}
We started with two standalone queries. One counted titles beginning with
\texttt{Champion} or \texttt{Loser}. The other counted rows where
\texttt{info\_type\_id = 101} and the stored rating text satisfied
\texttt{info > '5.0'}. The scans gave the planner the following counts.

\begin{tabularx}{\linewidth}{@{}Xrr@{}}
\toprule
\textbf{Filtered input} & \textbf{Original estimate} & \textbf{Saved count} \\
\midrule
\texttt{title} & 470 & 294 \\
\texttt{movie\_info\_idx} & 169,381 & 352,990 \\
\bottomrule
\end{tabularx}

We then joined the inputs on \texttt{title.id = movie\_info\_idx.movie\_id}.
The planner used both scan counts to estimate the join. The join returned
49 rows, which the store saved under a separate key for the pair.

Finally, we planned a query that kept the same pair and added
\texttt{kind\_type} to select movies. The planner reused 49 as the estimate
for the pair and computed the final three-table estimate separately.
Adding \texttt{kind\_type.id = 1} lets the planner infer
\texttt{title.kind\_id = 1}. That adds a scan filter, so the old pair count
no longer matches. The benchmark's \texttt{DEBUG1} messages show two scan
hits and one pair hit for the matching query. The changed query reports no
pair hit.

\newpage

\section*{Verification and results}
We ran the same IMDb queries on stock PostgreSQL 16.13 and our patched server
with feedback enabled. Both used the same data, indexes and server settings.
The timed queries had a 15-second limit and three measurements per query.
The separate screening queries had a four-second limit.

\begin{tabularx}{\linewidth}{@{}Xrr@{}}
\toprule
\textbf{IMDb benchmark} & \textbf{Stock PostgreSQL} & \textbf{Our project} \\
\midrule
Timed queries completed out of 28 & 21 & 21 \\
Screening queries completed out of 85 & 57 & 67 \\
Median Q-error for the same 19 timed queries & 53 & 1 \\
Geometric mean speedup for those 19 queries & \(1.00\times\) & \(1.76\times\) \\
\bottomrule
\end{tabularx}

A Q-error of 1 is the best score. The 19 queries in the last two rows
completed in both setups. Our project ran faster on 14 and slower on 5 of
those queries. Binary results agreed whenever both setups returned a result
that could be compared.

\begin{tabularx}{\linewidth}{@{}Xr@{}}
\toprule
\textbf{Checks on the assertion-enabled build} & \textbf{Passed} \\
\midrule
PostgreSQL regression tests & 219 \\
PostgreSQL isolation tests & 117 \\
Project checks for reuse, partial execution and concurrency & 46 \\
\bottomrule
\end{tabularx}

Run \texttt{bash project/run.sh build} and then
\texttt{bash project/run.sh verify} to run the project checks. The benchmark
instructions and saved results are under \texttt{results/job/}.

\section*{Contributions}
\textbf{Prabhav Khandelwal} worked on the proposal and checkpoint, including
the first PostgreSQL build, row counting and the initial feedback store.

\textbf{Sanchita Chaurasia} worked on join feedback, integration checks and
reuse of scan and two-table join counts in later queries.

\textbf{Saksham Saharia} worked on the IMDb setup, benchmark runs and
comparison of query plans and results and also implemented support for
cardinality-feedback history, confidence diagnostics and stale-entry pruning.

\textbf{AI assistance.} AI implemented the snapshot checks and reuse keys
and helped automate benchmarks and tests.

\begin{thebibliography}{1}
\bibitem{leo} M. Stillger, G. Lohman, V. Markl, and M. Kandil.
\emph{LEO -- DB2's LEarning Optimizer}. VLDB, 2001.
\end{thebibliography}

\end{document}
